% generated by roam.py handbook export -- не редактировать, править заметки
% 20260922182947-оценка_для_выпуклых_функции_с_липшицевым_градиентом.org (body 8829a1d674d5)

\begin{theorem}[{Оценка для выпуклых функций с липшицевым градиентом}]\label[theorem]{note:655064DF-D686-4726-83DC-6F3D67CC5B9C}
  Пусть $f: \R^n \rightarrow \R$, $f \in C^{1,1}_L$ (\hyperref[note:7DD1C7AB-1A49-4675-9F21-5A16C717F670]{липшицев градиент}) и \hyperref[note:1CA14F2C-637F-41DE-9E97-72075D19E97F]{выпуклая};
  Тогда
  \[
  \forall x, y \in \R^n\ f(x) + \left\langle \nabla f(x); y-x \right\rangle + \frac{1}{2L} \| \nabla f(y) - \nabla f(x) \|^2 \le f(y)
  \]
\end{theorem}

\begin{proof}
  Давайте рассмотрим $\phi(y) := f(y) + \left\langle - \nabla f(x); y \right\rangle$.
  Заметим, что $\phi$ выпуклая, так как $f$~--- выпуклая, а скалярное произведение~--- линейная функция;
  Тогда $\phi$~--- выпуклая, как сумма выпуклой и линейной функций

  Заметим, что у $\phi$ тоже липшицев градиент.
  \begin{multline*}
  \nabla \phi (y) = \nabla f (y) - \nabla f(x) \Rightarrow \| \nabla \phi(y_1) - \nabla \phi(y_2) \| = \\
  \| \nabla f (y_1) - \nabla f(x) - \nabla f (y_2) + \nabla f(x) \| = \| \nabla f(y_1) - \nabla f(y_2) \|
  \end{multline*}
  А мы знаем, что у $f$ липшицев градиент. Тогда
  \[
  \| \nabla \phi(y_1) - \nabla \phi(y_2) \| = \| \nabla f(y_1) - \nabla f(y_2) \| \le L \|y_1 - y_2 \|
  \]

  Заметим, что $\nabla \phi(x) = \nabla f(x) - \nabla f(x) = 0$. $\phi$ выпуклая, поэтому по \hyperref[note:035B8935-65A9-4BB4-BDA3-60A60FF278F0]{теореме о нулевом градиенте} $x$~--- точка наименьшего значения $\phi$.

  Тогда $\phi(x) \le \phi(z)\ \forall z \in \R^n$.
  Возьмем $z = y - \frac{1}{L} \nabla \phi(y)$.

  По \hyperref[note:379E4CAD-A44B-4631-8DAF-F3DC708EAEDE]{усиленному неравенству} для $\phi$ (у нее липшицев градиент) $\phi(z)-\phi(y) - \left\langle \nabla \phi(y);z - y \right\rangle \le \frac{L}{2} \| z -y \|^2$ \\
  Давайте оценим $\phi(z)$. И подставим $z$ в правой части
  \[
  \phi(z) \le \phi(y) +  \left\langle \nabla \phi (y); y - \frac{1}{L} \nabla \phi (y) - y \right\rangle + \frac{L}{2} \left\| y - \frac{1}{L} \nabla \phi (y) - y \right\|^{2}
  \]
  Сократим
  \[
  \phi (z) \le  \phi(y) - \frac{1}{L} \left\langle \nabla \phi(y); \nabla \phi(y) \right\rangle + \frac{1}{2L} \| \nabla \phi(y) \|^2 =
  \phi(y) - \frac{1}{2L} \| \nabla \phi (y) \|^2
  \]

  Получили $\phi(x) \le \phi(y) - \frac{1}{2L} \| \nabla \phi (y) \|^2$, подставим значение $\phi$
  \[
  f(x) - \left\langle \nabla f(x); x \right\rangle \le f(y) - \left\langle \nabla f(x);y \right\rangle - \frac{1}{2L} \| \nabla f(y) - \nabla f(x) \|^2
  \]
  Если оставить справа только $f(y)$ и сократить скалярные, то получим
  \[
  f(x) + \left\langle \nabla f(x); y -x \right\rangle + \frac{1}{2L} \| \nabla f(y) - \nabla f(x) \|^2  \le f(y)
  \]
\end{proof}
